% ===========================================================================
%  Hw02 Redone
%
%  THE WRITE-UP, and it is the assistant's to type.
%    The assistant transcribes each question or topic faithfully and emits an
%    empty, marked solution region beneath it. The student does the
%    mathematics; the assistant typesets what they agreed was right.
% ===========================================================================
\documentclass[11pt]{article}
\usepackage{coursemacros}

\title{Homework 02 --- Problems 61 and 76, Redone}
\author{Mikey Ferguson}
\date{\today}

\begin{document}
\maketitle

% The two problems the eleventh edition replaced, redone. The statements below
% are the eleventh edition's, read off textbook/ProbTextbook.pdf. The tenth
% edition's versions of these same two numbers are in
% homework/hw02/superseded-edition-10.tex and are not what was assigned.
% ---------------------------------------------------------------------------

\begin{problem}{61}
  (Edition 11, Chapter 2.) Let $X_1, X_2, \ldots$ be a sequence of independent,
  identically distributed continuous random variables. We say that a
  \emph{record} occurs at time $n$ if $X_n > \max(X_1, \ldots, X_{n-1})$ ---
  that is, if $X_n$ is larger than each of $X_1, \ldots, X_{n-1}$. Show that
  \begin{enumerate}
    \item $P\{\text{a record occurs at time } n\} = 1/n$;
    \item $\E[\text{number of records by time } n] = \sum_{i=1}^{n} 1/i$;
    \item $\Var(\text{number of records by time } n) = \sum_{i=1}^{n} (i-1)/i^2$;
    \item letting $N = \min\{n : n > 1 \text{ and a record occurs at time } n\}$,
          $\E[N] = \infty$.
  \end{enumerate}
  \emph{Hint for (b) and (c): represent the number of records as a sum of
  indicator variables.}
\end{problem}
% ===== SOLUTION 61 =====
\paragraph{(a) $P\{\text{a record occurs at time } n\} = 1/n$.}
The $X_i$ are continuous, so $P\{X_i = X_j\} = 0$ for $i \neq j$ and with
probability $1$ the $n$ values are distinct and can be ranked. Because
$X_1, \ldots, X_n$ are independent and identically distributed, relabelling
them leaves the joint distribution unchanged, so each of the $n!$ orderings of
$X_1, \ldots, X_n$ from smallest to largest carries the same probability
$1/n!$.

A record at time $n$ says exactly one thing about the ordering: $X_n$ is larger
than every other value, so it sits last. Fixing $X_n$ in the last place leaves
the remaining $n-1$ values to arrange freely, which happens in $(n-1)!$ ways.
Hence
\[
  P\{X_n > \max(X_1, \ldots, X_{n-1})\}
    = \frac{(n-1)!}{n!}
    = \frac{1}{n}.
\]

\paragraph{(b) $\E[\text{number of records by time } n] = \sum_{i=1}^{n} 1/i$.}
Let
\[
  I_i =
  \begin{cases}
    1 & \text{if a record occurs at time } i,\\
    0 & \text{otherwise,}
  \end{cases}
\]
so that $I_i$ is a random variable taking the values $0$ and $1$, and the number of
records by time $n$ is the sum $N = \sum_{i=1}^{n} I_i$. Each record contributes exactly
$1$ to that sum and each non-record contributes $0$, so the sum counts them.
\begin{align*}
  \E[\text{number of records by time } n]
    &= \EE{\sum_{i=1}^{n} I_i}\\
    &= \sum_{i=1}^{n} \EE{I_i}\\
    &= \sum_{i=1}^{n} \left(\frac{1}{i}\cdot 1 + \frac{i-1}{i}\cdot 0\right)\\
    &= \sum_{i=1}^{n} \frac{1}{i} .
\end{align*}

\paragraph{(c) $\Var(\text{number of records by time } n) = \sum_{i=1}^{n} (i-1)/i^2$.}
Keep the indicators $I_i$ and the identity $N = \sum_{i=1}^{n} I_i$ from part (b).
\begin{align*}
  \Var(N)
    &= \Var\!\left(\sum_{i=1}^{n} I_i\right)\\
    &= \sum_{i=1}^{n} \Var(I_i)\\
    &= \sum_{i=1}^{n} \left(\EE{I_i^2} - \EE{I_i}^2\right)\\
    &= \sum_{i=1}^{n}
       \left(1^2\cdot\frac{1}{i} + 0^2\cdot\frac{i-1}{i}
             - \left(1\cdot\frac{1}{i} + 0\cdot\frac{i-1}{i}\right)^{2}\right)\\
    &= \sum_{i=1}^{n} \left(\frac{1}{i} - \frac{1}{i^2}\right)
     = \sum_{i=1}^{n} \frac{i-1}{i^2} .
\end{align*}
The second line is additivity of variance over a sum due to independence.

\paragraph{(d) $N = \min\{n : n > 1 \text{ and a record occurs at time } n\}$ has
$\E[N] = \infty$.}
$N$ takes values in $\{2, 3, \ldots\}$, since time $1$ is a record by convention and
$N$ is the first record time after it. So
\[
  \E[N] = \sum_{i=2}^{\infty} i\, P\{N = i\} .
\]
Fix $i \ge 2$ and count orderings of $X_1, \ldots, X_i$, of which all $i!$ are equally
likely. The event $\{N = i\}$ says two things: no record occurred at any of the times
$2, \ldots, i-1$, which means nothing among $X_2, \ldots, X_{i-1}$ ever overtook $X_1$,
so $X_1$ is the largest of $X_1, \ldots, X_{i-1}$; and a record occurred at time $i$,
so $X_i$ is the largest of all $i$. Those two constraints pin the last two places in the
ordering --- $X_i$ largest, $X_1$ next --- and leave the remaining $i-2$ values to be
arranged freely, in $(i-2)!$ ways. Hence
\[
  P\{N = i\} = \frac{(i-2)!}{i!} .
\]
Substituting, and using $i! = i\,(i-1)!$ to cancel the leading $i$,
\[
  \E[N]
    = \sum_{i=2}^{\infty} i\,\frac{(i-2)!}{i!}
    = \sum_{i=2}^{\infty} \frac{(i-2)!}{(i-1)!}
    = \sum_{i=2}^{\infty} \frac{1}{i-1}
    = \sum_{m=1}^{\infty} \frac{1}{m} ,
\]
the harmonic series, which diverges. So $\E[N] = \infty$.

% ===== END SOLUTION 61 =====

\begin{problem}{76}
  (Edition 11, Chapter 2.) Use Chebyshev's inequality to prove the weak law of
  large numbers. Namely, if $X_1, X_2, \ldots$ are independent and identically
  distributed with mean $\mu$ and variance $\sigma^2$, then for any
  $\varepsilon > 0$,
  \[
    P\left\{ \left| \frac{X_1 + X_2 + \cdots + X_n}{n} - \mu \right|
             > \varepsilon \right\} \to 0
    \qquad \text{as } n \to \infty .
  \]
\end{problem}
% ===== SOLUTION 76 =====
Write $\bar{X}_{n} = \dfrac{X_{1} + X_{2} + \cdots + X_{n}}{n}$ for the sample mean.
Chebyshev's inequality says that for a random variable $Y$ with finite mean and variance
and any $k > 0$,
\[
  \PR{\,\left| Y - \EE{Y} \right| \ge k\,} \le \frac{\Var(Y)}{k^{2}} ,
\]
so the two things needed before it can be applied to $\bar{X}_{n}$ are $\EE{\bar{X}_{n}}$
and $\Var(\bar{X}_{n})$.

The mean is $\mu$. By linearity of expectation,
\[
  \EE{\bar{X}_{n}}
    = \frac{1}{n}\,\EE{X_{1} + \cdots + X_{n}}
    = \frac{1}{n}\left( \EE{X_{1}} + \cdots + \EE{X_{n}} \right)
    = \frac{n\mu}{n}
    = \mu ,
\]
so the quantity inside the absolute value in the statement is exactly
$\bar{X}_{n} - \EE{\bar{X}_{n}}$, which is the form Chebyshev requires.

The variance shrinks like $1/n$. The constant $1/n$ leaves a variance squared, and because
the $X_{i}$ are independent the variance of the sum is the sum of the variances:
\[
  \Var(\bar{X}_{n})
    = \frac{1}{n^{2}} \sum_{i=1}^{n} \Var(X_{i})
    = \frac{n\sigma^{2}}{n^{2}}
    = \frac{\sigma^{2}}{n} .
\]

Now fix $\varepsilon > 0$ and apply Chebyshev to $Y = \bar{X}_{n}$ with $k = \varepsilon$.
The event in the problem is the strict one, $\left| \bar{X}_{n} - \mu \right| >
\varepsilon$, and it is contained in the non-strict event $\left| \bar{X}_{n} - \mu \right|
\ge \varepsilon$ that Chebyshev bounds, so the bound applies to it as well:
\[
  \PR{\,\left| \bar{X}_{n} - \mu \right| > \varepsilon\,}
    \le \PR{\,\left| \bar{X}_{n} - \mu \right| \ge \varepsilon\,}
    \le \frac{\Var(\bar{X}_{n})}{\varepsilon^{2}}
    = \frac{\sigma^{2}}{n\varepsilon^{2}} .
\]
With $\sigma^{2}$ and $\varepsilon$ both fixed, the right-hand side is a constant divided
by $n$, so it tends to $0$ as $n \to \infty$ whatever the value of $\varepsilon$. A
probability is non-negative, so the left-hand side is squeezed to $0$ as well:
\[
  \PR{\,\left| \frac{X_{1} + X_{2} + \cdots + X_{n}}{n} - \mu \right| > \varepsilon\,}
    \longrightarrow 0
  \qquad \text{as } n \to \infty ,
\]
% ===== END SOLUTION 76 =====

\end{document}
